% 2019 数学一 q06：三张平面两两相交，三条交线互相平行
% 可实现的空间模型：x=0, y=0, x+y=1。
% 三平面沿 z 方向延展；两两交线分别为 (0,0,z), (0,1,z), (1,0,z)，严格平行且无公共点。

  % 用斜投影画出底面三角形的三个顶点，以及沿共同 z 方向平移后的对应点
  \coordinate (A)  at (0.0, 0.0);
  \coordinate (B)  at (3.6, 0.25);
  \coordinate (C)  at (1.35, 1.55);
  \coordinate (A2) at (0.55, 3.20);
  \coordinate (B2) at (4.15, 3.45);
  \coordinate (C2) at (1.90, 4.75);

  % 三张平面的有限补丁：分别由底面三角形的三条边沿共同方向拉伸得到
  % pi_1 : x=0，对应 AC 边
  \draw[line width=0.9pt] (A) -- (C) -- (C2) -- (A2) -- cycle;
  \node at (0.70, 2.25) {$\pi_1$};

  % pi_2 : y=0，对应 AB 边
  \draw[line width=0.9pt] (A) -- (B) -- (B2) -- (A2) -- cycle;
  \node at (2.55, 1.85) {$\pi_2$};

  % pi_3 : x+y=1，对应 BC 边
  \draw[line width=0.9pt] (B) -- (C) -- (C2) -- (B2) -- cycle;
  \node at (3.18, 2.75) {$\pi_3$};

  % 三条两两交线：AA2、BB2、CC2，共享完全相同的方向向量 (0.55, 3.20)
  \draw[line width=1.35pt] (A) -- (A2);
  \draw[line width=1.35pt] (B) -- (B2);
  \draw[line width=1.35pt] (C) -- (C2);

  % 用同向箭头强调三条交线平行
  \draw[->, line width=0.7pt] (0.20,1.15) -- (0.34,1.95);
  \draw[->, line width=0.7pt] (3.80,1.35) -- (3.94,2.15);
  \draw[->, line width=0.7pt] (1.55,2.45) -- (1.69,3.25);
